% Pursuing Stacks (cotes 134-2 to 134-8): archive cote 134-8, site page 13 (PDF page 14), typescript page 567.
% typescript. Draft compared: ps7.tex:600-659.
% AI-transcribed by Claude Opus 5.5 (claude-opus-5-5, adaptive thinking, effort high), Message Batches API, batch ps4 (msgbatch_012Aoc1NuskATPsyBeKSYK7J), 2026-11-12; not yet checked by eye.
\scanpage{134-8/13}
\mspage{567}

from the former,  using the relevant two functors between (Cat) and \hand{$\Delta$}$^{\wedge}$
which allow to pass from one type of models to another. (Soon or late
I should check in Bousf\unsure{ie}ld-Kan's book\xout{s} whether this is so or not ...)
Thomason however did the opposite,  and it is quite natural that he had
to pay for it by a fair amount of sweat!  \xout{xx} (Reference of the paper:
Homotopy colimits in the category of small categories,  Math. Proc.
Camb. Phil. Soc. (1979)\unsure{;} 85, 91,  p. 91 - 109.)

The proof written down yesterday for compatibility of abelianizati-%
on of homotopy types with the Whitehead abelianization functor within
a category A$^{\wedge}$,  still goes through when replacing abelianization by
k-linearization,  with respect to an arbitrary ground ring k (not even
commutative).  It was'nt really worth while though to introduce a ring
k,  as the general result should follow at once from the case k=$\underline{Z}$ ,
by the formula
\[ (1) \qquad Lp^{k}_{!}(k_{A}) \simeq Lp^{\underline{Z}}_{!}(\underline{Z}_{A}) \overset{L}{\otimes}_{\underline{Z}} k \quad , \]
where in the \xout{first}\inserted{left} hand side we are taking the \inserted{left} derived functor for the
functor
\[ p^{k}_{!} \ : \ A^{\wedge}_{k} \longrightarrow (Ab_{k}) \]
generalizing $p^{ab}_{!} = p^{\underline{Z}}_{!}$ , \hand{(}\over{where}{with} $p\colon A \to e$ \deleted{is} as above\hand{)}, and where in
the \xout{second member} right hand side we are using the \unsure{ring} extension
functor
\[ \overset{L}{\otimes}_{\underline{Z}} \ : \ D^{\bullet}(Ab) \longrightarrow D_{.}(Ab_{k}) \]
for the relevant derived categories.  This reminds me of the need of
developping \xout{more} a more or less exhaustive formulaire around the
basic operations
\[ (2) \qquad Lf_{!} \ , \ f^{*} \ , \ Rf_{*} \quad , \]
including the familiar one for  the two latter,  valid \xout{for} more general-%
ly for maps between ringed topoi,  and including also \inserted{a} \hand{``}\uline{projection
formula}'' generalizing (1).  Such a formula will be no surprise surely
to a reader familiar with a duality context (such as étale cohomology,
of ``coherent'' cohomology of noetherian schemes say )\hand{,} where a formalism
of the ``four variances''  $f_{!}$ , $f^{*}$ , $f_{*}$ , $f^{!}$  and the \hand{``}two internal opera-%
tions\hand{''} $\overset{L}{\otimes}$  and R\uline{Hom}  can be developped -  it would seem that the formal
properties of the triple (2),  together with the two internal operations
just referred to,  are very close (for the least) \over{\unsure{of}}{to} those of the
slightly richer one,\inserted{in duality set-ups,} including equally an ``unusual inverse image'' $f^{!}$,
\margin{no, to $Rf_!$.} right adjoint to \uline{$Rf_{*}$}  (denoted sometimes simply by $f_{*}$ ).  This similari

\margin{13}
